\documentclass[a4paper,10pt]{article}

% --- PAQUETES DE DISEÑO ---
\usepackage[utf8]{inputenc}
\usepackage[T1]{fontenc}
\usepackage[default]{lato} % Fuente moderna sans-serif
\usepackage[margin=1.25cm]{geometry} % Márgenes aprovechados
\usepackage{xcolor}
\usepackage{hyperref}
\usepackage{fontawesome5} % Para los iconos (LinkedIn, Mail...)
\usepackage{titlesec}
\usepackage{enumitem}

% --- COLORES (Personalizables) ---
\definecolor{primary}{RGB}{40, 100, 160} % Azul oscuro profesional
\definecolor{secondary}{RGB}{80, 80, 80} % Gris oscuro

% --- ESTILO DE TÍTULOS ---
% Título de sección con línea debajo
\titleformat{\section}
{\Large\bfseries\color{primary}}
{}{0em}
{}[\titlerule]
\titlespacing*{\section}{0pt}{1.2em}{0.6em}

% --- QUITAR NUMERACIÓN ---
\pagestyle{empty}

% --- COMIENZO DEL DOCUMENTO ---
\begin{document}

% === CABECERA ===
\begin{center}
    {\Huge \textbf{Tu Nombre}} \\ \vspace{0.2cm}
    {\Large \textcolor{secondary}{}} \\ \vspace{0.2cm}
    
    % Iconos de contacto (usamos condicionales por si algún dato falta)
    \small
                \end{center}

\vspace{0.3cm}

% === PERFIL ===

% === EXPERIENCIA ===
\section*{Experiencia Laboral}
\begin{itemize}[leftmargin=0cm, label={}]
        \item 
        \textbf{\large Investigador PhD} \hfill \textbf{} \\
        \textcolor{primary}{\textit{Universidad}} \ \textbullet \ \textit{} \\
        \vspace{0.05cm}
        Análisis de datos clínicos y redacción de papers.
        \vspace{0.3cm}
    \end{itemize}

% === FORMACIÓN ===
% Solo mostramos la sección si hay datos en la lista filtrada

% === HABILIDADES ===

\end{document}